Die automatisierte, referenzfreie Bewertung von Sprachkomplexität stellt eine Kernkomponente dieser Arbeit dar. Um Textvereinfachungen objektiv zu bewerten und ein steuerbares Belohnungssignal (\textit{Reward}) für das spätere Modelltraining bereitzustellen, wird in diesem Themenblock eine eigene Metrik-Pipeline entwickelt und systematisch untersucht.

\subsection{Klassifikationsmodelle (BiLSTM vs. Baselines)}
\label{sec:klassifikationsmodelle}
\begin{tcolorbox}[colback=red!5!white,colframe=red!75!black,title=\textbf{Zu erklären in diesem Abschnitt:},sharp corners,boxrule=0.8pt]
\color{red!80!black}
\begin{itemize}
    \item \textbf{[Zu erklären] Modellarchitektur:} Bidirektionales Long Short-Term Memory (BiLSTM) mit Word-Embedding-Layer, Dropout-Regularisierung und Klassifikationskopf (Linear + Sigmoid).
    \item \textbf{[Zu erklären] Vergleichsbaselines:} Vanilla RNN, Standard-LSTM und vortrainierte Transformer-Klassifikatoren (GBERT).
    \item \textbf{[Zu erklären] Granularitätsebenen:}
    \begin{itemize}
        \item \textit{Satz-Klassifikator:} Training und Inferenz auf isolierten Einzelsätzen.
        \item \textit{Artikel-Klassifikator:} Dokumentenweites Training (bis 512 Tokens).
        \item \textit{Majority Voting:} Aggregierte Dokumenten-Klassifikation durch Mehrheitsentscheid über alle Sätze eines Textes.
    \end{itemize}
    \item \textbf{[Zu erklären] In-Domain-Evaluation:} Leistungsvergleich anhand von Balanced Accuracy (BAcc) und F1-Score auf dem internen Test-Split.
\end{itemize}
\end{tcolorbox}

\subsection{Out-of-Domain-Generalisierung \& Empirische Bias-Kontrolle}
\label{sec:ood_bias_kontrolle}
\begin{tcolorbox}[colback=red!5!white,colframe=red!75!black,title=\textbf{Zu erklären in diesem Abschnitt:},sharp corners,boxrule=0.8pt]
\color{red!80!black}
\begin{itemize}
    \item \textbf{[Zu erklären] Out-of-Domain-Validierung:} Evaluation auf einem unabhängigen, externen Testdatensatz der \textit{Lebenshilfe Schleswig-Holstein e.\,V.}
    \item \textbf{[Zu erklären] Shortcut- und Bias-Kontrollen:}
    \begin{itemize}
        \item \textit{Dummy-Content-Test (Längen-Bias):} Maskierung aller Wörter durch identische Satzzeichen (`\texttt{.}') bei unveränderter Tokenanzahl zur Kontrolle von Längen-Shortcuts.
        \item \textit{Layout-Bias-Test:} Vollständige Entfernung von Absätzen und Zeilenumbrüchen zur Überprüfung von typografischen Heuristiken.
    \end{itemize}
\end{itemize}
\end{tcolorbox}

\subsection{Continuous MixUp Regression}
\label{sec:mixup_regression}
\begin{tcolorbox}[colback=red!5!white,colframe=red!75!black,title=\textbf{Zu erklären in diesem Abschnitt:},sharp corners,boxrule=0.8pt]
\color{red!80!black}
\begin{itemize}
    \item \textbf{[Zu erklären] Das Konzept des textuellen MixUp:} Erzeugung kontinuierlicher Schwierigkeitsgrade $\lambda \in [0{,}0, 1{,}0]$ durch satzweise Mischungen paralleler Artikel.
    \item \textbf{[Zu erklären] Dynamische Target-Berechnung:} Gewichtung anhand des relativen Zeichenlängenanteils:
    \begin{equation}
        \lambda = \frac{\text{CharLen}(LS)}{\text{CharLen}(LS) + \text{CharLen}(AS)}
    \end{equation}
    \item \textbf{[Zu erklären] Vergleich der Trainings- und Dataloader-Varianten:}
    \begin{itemize}
        \item \textit{Variante A (Statisch):} Feste Vorabmischung des Trainingskorpus.
        \item \textit{Variante B (Dynamisch):} On-the-Fly Shuffling (Analyse des \textit{Prediction Collapse}).
        \item \textit{Variante C (Hybrid):} Epochenabhängiges Curriculum-Shuffling.
        \item \textit{Variante D (Hybrid + Cyclic Learning Rate):} Optimierung mittels \textit{CosineAnnealingWarmRestarts} zur Vermeidung lokaler Minima.
    \end{itemize}
    \item \textbf{[Zu erklären] Kontextlängen-Ablationsstudie:} Empirische Bestimmung der optimalen Eingabesequenzlänge (128, 256, 500, 1000 Tokens).
    \item \textbf{[Zu erklären] Regressions-Metriken:} Mean Squared Error (MSE), Mean Absolute Error (MAE), Pearson- und Spearman-Korrelation sowie KDE-Dichteverteilungsanalysen.
\end{itemize}
\end{tcolorbox}

\subsection{Evaluierung synthetischer Sprachstufen}
\label{sec:synthetische_sprachstufen}
\begin{tcolorbox}[colback=red!5!white,colframe=red!75!black,title=\textbf{Zu erklären in diesem Abschnitt:},sharp corners,boxrule=0.8pt]
\color{red!80!black}
\begin{itemize}
    \item \textbf{[Zu erklären] LLM-basierte Stufengenerierung:} Synthetische Erzeugung diskreter Zwischenstufen ($0{,}25$, $0{,}50$, $0{,}75$) mittels Instruction Prompting großer Sprachmodelle (FlensGen-GPT / OpenAI).
    \item \textbf{[Zu erklären] Kreuz-Evaluation:} Gegenüberstellung des MixUp-Regressors und des synthetisch trainierten Regressors auf dem Lebenshilfe-Datensatz.
    \item \textbf{[Zu erklären] Monotonie- und Reliabilitätsprüfung:} Analyse von Fehlertrends, Boxplots und stilistischen Artefakten (\textit{Chatty Prefixes}, Formatierungsverluste).
\end{itemize}
\end{tcolorbox}
